% Standalone compilation: run pdflatex abstract.tex.
% In the main manuscript, define \FINELManuscriptMode before % Standalone compilation: run pdflatex abstract.tex.
% In the main manuscript, define \FINELManuscriptMode before % Standalone compilation: run pdflatex abstract.tex.
% In the main manuscript, define \FINELManuscriptMode before % Standalone compilation: run pdflatex abstract.tex.
% In the main manuscript, define \FINELManuscriptMode before \input{abstract.tex}.
\ifdefined\FINELManuscriptMode
\else
\documentclass[a4paper,11pt]{article}
\usepackage[T1]{fontenc}
\usepackage[a4paper,margin=1in]{geometry}
\usepackage{newtxtext,newtxmath}
\begin{document}
\fi
\begin{abstract}
The open-circuit voltage of a tetragonal BaTiO$_3$ microcantilever under a steady through-thickness temperature difference is examined using a coupled thermopiezoelectric finite-element model. The plane-strain formulation follows from electric enthalpy and employs the fixed-strain pyroelectric coefficient. The thermal-stress vector is obtained by contracting the full three-dimensional stiffness and thermal-expansion tensors before imposing plane strain. Omitting the constrained transverse expansion underestimates two thermal-stress components by 35\% and 58\%, respectively, and reverses the predicted piezoelectric--thermal voltage from $+0.19$~V to $-0.22$~V. The voltage is $-1.50$~V when pyroelectric and piezoelectric couplings are both active. The coupled result cannot be reconstructed by adding the separately computed limiting-case voltages because electromechanical coupling also alters the pyroelectric contribution. The formulation is checked against an exact one-dimensional pyroelectric limit, an analytical traction-free interior solution, mesh refinement and an independent biquadratic discretization. Effects of the root constraint and uncertainty in the fitted thermal-expansion coefficients are also assessed. The results show how neglecting blocked transverse thermal expansion can lead to an incorrect voltage polarity in plane-strain calculations of anisotropic piezoelectric devices.
\end{abstract}
\ifdefined\FINELManuscriptMode
\else
\end{document}
\fi
.
\ifdefined\FINELManuscriptMode
\else
\documentclass[a4paper,11pt]{article}
\usepackage[T1]{fontenc}
\usepackage[a4paper,margin=1in]{geometry}
\usepackage{newtxtext,newtxmath}
\begin{document}
\fi
\begin{abstract}
The open-circuit voltage of a tetragonal BaTiO$_3$ microcantilever under a steady through-thickness temperature difference is examined using a coupled thermopiezoelectric finite-element model. The plane-strain formulation follows from electric enthalpy and employs the fixed-strain pyroelectric coefficient. The thermal-stress vector is obtained by contracting the full three-dimensional stiffness and thermal-expansion tensors before imposing plane strain. Omitting the constrained transverse expansion underestimates two thermal-stress components by 35\% and 58\%, respectively, and reverses the predicted piezoelectric--thermal voltage from $+0.19$~V to $-0.22$~V. The voltage is $-1.50$~V when pyroelectric and piezoelectric couplings are both active. The coupled result cannot be reconstructed by adding the separately computed limiting-case voltages because electromechanical coupling also alters the pyroelectric contribution. The formulation is checked against an exact one-dimensional pyroelectric limit, an analytical traction-free interior solution, mesh refinement and an independent biquadratic discretization. Effects of the root constraint and uncertainty in the fitted thermal-expansion coefficients are also assessed. The results show how neglecting blocked transverse thermal expansion can lead to an incorrect voltage polarity in plane-strain calculations of anisotropic piezoelectric devices.
\end{abstract}
\ifdefined\FINELManuscriptMode
\else
\end{document}
\fi
.
\ifdefined\FINELManuscriptMode
\else
\documentclass[a4paper,11pt]{article}
\usepackage[T1]{fontenc}
\usepackage[a4paper,margin=1in]{geometry}
\usepackage{newtxtext,newtxmath}
\begin{document}
\fi
\begin{abstract}
The open-circuit voltage of a tetragonal BaTiO$_3$ microcantilever under a steady through-thickness temperature difference is examined using a coupled thermopiezoelectric finite-element model. The plane-strain formulation follows from electric enthalpy and employs the fixed-strain pyroelectric coefficient. The thermal-stress vector is obtained by contracting the full three-dimensional stiffness and thermal-expansion tensors before imposing plane strain. Omitting the constrained transverse expansion underestimates two thermal-stress components by 35\% and 58\%, respectively, and reverses the predicted piezoelectric--thermal voltage from $+0.19$~V to $-0.22$~V. The voltage is $-1.50$~V when pyroelectric and piezoelectric couplings are both active. The coupled result cannot be reconstructed by adding the separately computed limiting-case voltages because electromechanical coupling also alters the pyroelectric contribution. The formulation is checked against an exact one-dimensional pyroelectric limit, an analytical traction-free interior solution, mesh refinement and an independent biquadratic discretization. Effects of the root constraint and uncertainty in the fitted thermal-expansion coefficients are also assessed. The results show how neglecting blocked transverse thermal expansion can lead to an incorrect voltage polarity in plane-strain calculations of anisotropic piezoelectric devices.
\end{abstract}
\ifdefined\FINELManuscriptMode
\else
\end{document}
\fi
.
\ifdefined\FINELManuscriptMode
\else
\documentclass[a4paper,11pt]{article}
\usepackage[T1]{fontenc}
\usepackage[a4paper,margin=1in]{geometry}
\usepackage{newtxtext,newtxmath}
\begin{document}
\fi
\begin{abstract}
The open-circuit voltage of a tetragonal BaTiO$_3$ microcantilever under a steady through-thickness temperature difference is examined using a coupled thermopiezoelectric finite-element model. The plane-strain formulation follows from electric enthalpy and employs the fixed-strain pyroelectric coefficient. The thermal-stress vector is obtained by contracting the full three-dimensional stiffness and thermal-expansion tensors before imposing plane strain. Omitting the constrained transverse expansion underestimates two thermal-stress components by 35\% and 58\%, respectively, and reverses the predicted piezoelectric--thermal voltage from $+0.19$~V to $-0.22$~V. The voltage is $-1.50$~V when pyroelectric and piezoelectric couplings are both active. The coupled result cannot be reconstructed by adding the separately computed limiting-case voltages because electromechanical coupling also alters the pyroelectric contribution. The formulation is checked against an exact one-dimensional pyroelectric limit, an analytical traction-free interior solution, mesh refinement and an independent biquadratic discretization. Effects of the root constraint and uncertainty in the fitted thermal-expansion coefficients are also assessed. The results show how neglecting blocked transverse thermal expansion can lead to an incorrect voltage polarity in plane-strain calculations of anisotropic piezoelectric devices.
\end{abstract}
\ifdefined\FINELManuscriptMode
\else
\end{document}
\fi
